% Generated draft: requires mathematical and source review.
\documentclass[11pt]{article}
\usepackage[margin=25mm]{geometry}
\usepackage{fontspec}
\setmainfont{Latin Modern Roman}
\usepackage{amsmath,amssymb,mathtools,unicode-math}
\setmathfont{Latin Modern Math}
\usepackage{enumitem,needspace,xurl,hyperref}
\setlength{\parindent}{0pt}
\setlength{\parskip}{5pt}
\setlength{\emergencystretch}{3em}
\newcommand{\N}{\mathbb{N}}
\newcommand{\Z}{\mathbb{Z}}
\newcommand{\Q}{\mathbb{Q}}
\newcommand{\R}{\mathbb{R}}
\newcommand{\C}{\mathbb{C}}
\newcommand{\F}{\mathbb{F}}
\newcommand{\A}{\mathbb{A}}
\newcommand{\PP}{\mathbb{P}}
\newcommand{\E}{\mathbb{E}}
\newcommand{\eps}{\varepsilon}
\newcommand{\dd}{\,\mathrm{d}}
\begin{document}
\section*{UnsolvedMath ID: 30005871}
\textbf{Source problem code:} catalogue entry
\subsection{Definitions and mathematical statement}
Let \(Y\) be an \(n\)-dimensional manifold with boundary \(\partial Y\), and write \(M=\operatorname{int}(Y)=Y\setminus\partial Y\). The supplied statement uses groups \(\Gamma'\) and \(\Gamma\), a homomorphism \(\iota\colon\Gamma'\to\Gamma\), and classifying data denoted by \(\nu\). For a discrete group \(G\), \(BG\) denotes a classifying space for \(G\). The intended classifying data must induce a homomorphism on relative homology carrying a rational relative fundamental class to
\[
\alpha=\nu_*[Y,\partial Y]\in H_n(B\Gamma,B\Gamma';\Q).
\]
Here \(H_n\) denotes relative singular homology with rational coefficients. The hypotheses ensuring the existence of \([Y,\partial Y]\), the precise classifying data, and the interpretation of the target pair are not specified in the supplied record. Thus the following questions are conditional on these data being well defined; they do not silently supply compactness, orientability, spin, or connectivity assumptions.

A positive-scalar-curvature metric on \(M\) means a smooth Riemannian metric \(g\) satisfying \(\operatorname{Scal}_g(x)>0\) for every \(x\in M\). Completeness means geodesic completeness. Finite volume means \(\operatorname{Vol}_g(M)<\infty\). No uniform positive lower bound for scalar curvature is stated.

\textbf{Intended obstruction questions.} Assuming \(\alpha\ne0\), does each of the following assertions hold separately?
\begin{enumerate}
\item If \(\iota\colon\Gamma'\to\Gamma\) is injective, then \(M\) admits no complete positive-scalar-curvature metric:
\[
\alpha\ne0\ \text{and}\ \iota\ \text{injective}
\quad\Longrightarrow\quad
\nexists\,g\text{ on }M\text{ with }g\text{ complete and }\operatorname{Scal}_g>0.
\]
\item Without additionally requiring injectivity, \(M\) admits no complete positive-scalar-curvature metric of finite volume:
\[
\alpha\ne0
\quad\Longrightarrow\quad
\nexists\,g\text{ on }M\text{ with }g\text{ complete},\
\operatorname{Scal}_g>0,\
\operatorname{Vol}_g(M)<\infty.
\]
\end{enumerate}
The finite-volume condition concerns the same metric whose completeness and scalar curvature are in question.

\textbf{Supplied special-case qualification.} The report excerpt identifies the case in which both \(Y\) and \(\partial Y\) are aspherical and asserts that \(\alpha\) is then automatically nonzero. For a connected manifold, aspherical means that its universal cover is contractible; a disconnected boundary would require a componentwise interpretation. The automatic nonvanishing is retained as a supplied source assertion, conditional on the intended fundamental-class and classifying-map setup. The excerpt does not provide enough surrounding hypotheses to state it here as an independently justified general fact.
\subsection{Short English statement}
The question asks whether a nonzero rational relative fundamental class, after pushforward to the indicated classifying spaces, rules out complete positive scalar curvature when the boundary-group homomorphism is injective, and separately rules out complete positive scalar curvature with finite volume. The source also highlights an aspherical special case. Missing hypotheses concerning the fundamental class and classifying data require review. The supplied status is \emph{partially solved}, not an independently verified conclusion.
\subsection{Sources}
\begin{itemize}
\item Local UnsolvedMath registry, permanent ID \(30005871\), title \emph{Relative Homology Obstructions to Complete Positive Scalar Curvature}, exact source problem code \texttt{catalogue entry}. The supplied record belongs to the Oberwolfach Reports Open Problems collection and assigns proposed year 2024; no individual proposer is identified.
\item \emph{Geometrie} (2024), Oberwolfach Reports, pp.\ 1563--1614, DOI \texttt{10.4171/owr/2024/28}: \url{https://doi.org/10.4171/owr/2024/28}. This is the authoritative report citation supplied by the record. No precise page or named contributor for this particular question is supplied. The registry labels the statement \texttt{corrected\_verified} and reports comparison with the official report; that verification is supplied metadata, not verification performed for this draft.
\item Supplied literature reference, ``Lott noncompact positive-scalar-curvature obstructions'': \url{https://doi.org/10.1515/crelle-2025-0072}.
\item Supplied literature reference, ``Lott primary preprint'': \url{https://arxiv.org/abs/2402.13239}.
\item Supplied literature reference, ``original OWR report'': \url{https://doi.org/10.4171/owr/2024/28}.
\item Literature assessment dated \texttt{2026-08-22}: the supplied curation labels the problem partially solved and says that Lott develops relative almost-flat and Dirac-index obstructions in the injective-boundary and finite-volume directions, presenting results toward the stated test question rather than a full theorem under either hypothesis. This is supplied dated curation, not an independently verified literature assessment; the record expressly requests independent verification of time-sensitive status and cited claims.
\end{itemize}
\end{document}
